\documentclass[english, parskip=full]{scrartcl}
\usepackage[utf8]{inputenc}  % Kodierung der Datei
\usepackage[english]{babel}  % Sprache auf Deutsch 
\usepackage{xcolor, colortbl}
\usepackage{graphicx, gensymb}
\usepackage{amsfonts, cancel, amssymb}
\usepackage{amsmath, amsthm, array, esint}
\usepackage{booktabs, multirow}  % schönere Tabellen
\usepackage{caption}  % Anpassung der Captions
\usepackage{scrlayer-scrpage}    % Manipulation des Seitenstils
\pagestyle{scrheadings}  % Stil für die Seite setzen
\automark{section}  % Abschnittsnamen als Seitenbeschriftung 
\setheadsepline{.5pt}    % Kopzeile durch Linie abtrennen
\usepackage[hidelinks]{hyperref}  % Links und weitere PDF-Features
\usepackage{typearea, tikz, pgffor, verbatim}
\usetikzlibrary{calc, arrows.meta,arrows, graphs, quotes}
\usepackage[cache=false, outputdir=build]{minted}
\usemintedstyle{vs}
\usepackage{placeins} % provides \FloatBarrier

\definecolor{bg}{rgb}{0.9,0.9,0.9}
\newcommand\numberthis{\addtocounter{equation}{1}\tag{\theequation}}
\usepackage{svg}
\usepackage{siunitx}
\usepackage{physics}
\usepackage{tensor}
\renewcommand{\bra}{}
\newcommand{\kom}{}
\newcommand{\brak}{}
\renewcommand{\braket}{}
\newcommand{\intx}{}
\newcommand{\intr}{}
\newcommand{\kla}{}
\newcommand{\klam}{}
\newcommand{\klammer}{}
\newcommand{\oper}{}
\newcommand{\tecxt}{}
\newcommand{\Bra}{}
\newcommand{\Ket}{}
\renewcommand{\i}{\text{i}}
\newcommand{\e}{\text{e}}
\newcommand*{\Fourier}{\textsc{Fourier}\ }
\newcommand*{\F}{\mathcal{F}}
\newcommand*{\sinc}{\text{sinc\,}}
\newcommand{\conv}{\mathop{\scalebox{3.5}{\raisebox{-0.38ex}{\makebox[0.5\width][c]{$\ast$}}}}\limits}

\newcommand*{\titel}{03 Geodesics and curvature}
\newcommand*{\autor}{Joachim Schwardt}

\hypersetup{pdfauthor={\autor}, pdftitle={\titel}}

%\titlehead{titlehead}
\subject{General Relativity}
\title{\titel}
\author{\autor}
\date{\today}

\begin{document}

%\begin{titlepage}
\maketitle  % Titel setzen
%\tableofcontents  % Inhaltsverzeichnis setzen
%\end{titlepage}


\section{Geodesics on a sphere}
For a sphere of radius $R$ we have $\tecxt\text{d}t^2 = R^2\tecxt\text{d}\Omega^2 = R^2 \kla\left( \tecxt\text{d}\theta^2 + \sin^2\theta \tecxt\text{d}\phi^2 \right)$.
Thus the metric is given by
\begin{align}
(g_{ab}) &= R^2 \begin{pmatrix}
1 & 0 \\ 0 & \sin^2\theta
\end{pmatrix}.
\end{align}
The Christoffel symbols are given by
\begin{align}
\Gamma^\theta_{ab} &= \frac{1}{2}g^{\theta c} \kla\left( g_{ca,b} + g_{cb,a} - g_{ab,c} \right) = \frac{1}{2R^2}\kla\left( g_{\theta a,\theta} \delta_{b\theta} + g_{\theta b,\theta} \delta_{a\theta} - g_{ab,\theta} \right) = \\
&= \frac{1}{2R^2} (-g_{\phi\phi,\theta}\delta_{a\phi} \delta_{b\phi}) = -\frac{1}{2}\delta_{a\phi} \delta_{b\phi}\partial_\theta \sin^2\theta = -\delta_{a\phi} \delta_{b\phi}\sin\theta\cos\theta ,\\
\Gamma^\phi_{ab} &= \frac{1}{2}g^{\phi c} \kla\left( g_{ca,b} + g_{cb,a} - g_{ab,c} \right) = \frac{1}{2R^2 \sin^2\theta}\kla\left( g_{\phi a,\theta} \delta_{b\theta} + g_{\phi b,\theta} \delta_{a\theta} - g_{ab,\phi} \right) = \\
&= \frac{1}{2R^2 \sin^2\theta} (g_{\phi\phi,\theta} \delta_{b\theta}\delta_{a\phi} + g_{\phi\phi,\theta} \delta_{a\theta} \delta_{b\phi}) = \cot\theta \kla\left( \delta_{b\theta} \delta_{a\phi} + \delta_{a\theta} \delta_{b\phi} \right).
\end{align}
Thus we find $\Gamma^\theta_{\phi\phi} = -\sin\theta\cos\theta$ and $\Gamma^\phi_{\theta\phi} = \cot\theta$.
\\
The Lagrangian is 
\begin{align}
L &= g_{ab}\dot{x}^a \dot{x}^b = R^2\kla\left( \dot{\theta}^2 + \sin^2\theta \dot{\phi}^2 \right).
\end{align}
The Euler-Lagrange equations are
\begin{align}
0 &= \frac{\tecxt\text{d}}{\tecxt\text{d}t}\frac{\partial L}{\partial \dot{\theta}} - \frac{\partial L}{\partial \theta} = \frac{\tecxt\text{d}}{\tecxt\text{d}t} \kla\left( 2R^2\dot{\theta} \right) - 2R^2\sin\theta\cos\theta\dot{\phi}^2 = \\
&= 2R^2\kla\left( \ddot{\theta} - \sin\theta\cos\theta\dot{\phi}^2 \right),\\
0 &= \frac{\tecxt\text{d}}{\tecxt\text{d}t}\frac{\partial L}{\partial \dot{\phi}} - \frac{\partial L}{\partial \phi} = \frac{\tecxt\text{d}}{\tecxt\text{d}t} \kla\left( 2R^2\sin^2\theta \dot{\phi} \right) = \\
&= 2R^2\sin^2\theta\kla\left( \ddot{\phi} + \frac{2\sin\theta\cos\theta}{\sin^2\theta} \dot{\theta}\dot{\phi} \right).
\end{align}
We can now read off the Christoffel symbols $\Gamma^\theta_{\phi\phi} = -\sin\theta\cos\theta$ and $\Gamma^\phi_{\theta\phi} = \cot\theta$, which is identical to the previous result.
\\
Now we can also figure out the motion on a sphere by solving the geodesics equation.
Here we have
\begin{align}
0 &= \ddot{\theta} + \Gamma^\theta_{\phi\phi}\dot{\phi}^2 = \ddot{\theta} - \sin\theta\cos\theta \dot{\phi}^2,\label{eqn:ddot theta} \\
0 &= \ddot{\phi} + 2\Gamma^\phi_{\theta\phi} \dot{\theta}\dot{\phi} = \ddot{\phi} + 2\cot\theta\dot{\theta}\dot{\phi}. \label{eqn:ddot phi}
\end{align}
\autoref{eqn:ddot phi} can be rewritten to
\begin{align}
\frac{\ddot{\phi}}{\dot{\phi}} &= -2\cot\theta\dot{\theta},\\
\Rightarrow\ \ \ \frac{\tecxt\text{d}}{\tecxt\text{d}t} \log\dot{\phi} &= -2\frac{\tecxt\text{d}}{\tecxt\text{d}t} \log \sin\theta,\\
\Rightarrow\ \ \ \dot{\phi} &= \exp\kla\left( -2\log\sin\theta + \tecxt\text{const} \right) = \frac{a}{\sin^2\theta},
\end{align}
where $a = \dot{\phi}_0 \sin^2\theta_0$ is determined by the initial conditions (this result is also obvious when considering the Euler-Lagrange equation, since $\phi$ is cyclic).
We may use this in \autoref{eqn:ddot theta} to get
\begin{align}
\ddot{\theta} \dot{\theta} &= \sin\theta\cos\theta \frac{a^2}{\sin^4\theta} \dot{\theta},\\
\Rightarrow\ \ \ \frac{\tecxt\text{d}}{\tecxt\text{d}t} \frac{\dot{\theta}^2}{2} &= \frac{\tecxt\text{d}}{\tecxt\text{d}t} \kla\left( -\frac{a^2}{2\sin^2\theta} \right),\\
\Rightarrow\ \ \ \dot{\theta} &= \sqrt{b^2 - \frac{a^2}{\sin^2\theta}},
\end{align}
where $b^2 = \dot{\theta}_0^2 + \dot{\phi}_0^2 \sin^2\theta_0$ is also determined by the initial conditions.
This equation is equivalent to
\begin{align}
bt + c &= \intx\int_{ }^{ } \frac{b\sin\theta}{\sqrt{b^2 \sin^2\theta - a^2}}\, \mathrm{d}\theta =: I(\theta).
\end{align}
We solve this integral by substituting $z = R\cos\theta$, which yields
\begin{align}
I(\theta) &= \intx\int_{ }^{ } \frac{Rb\,\tecxt\text{d}z}{\sqrt{b^2 - a^2 - b^2R^2z^2}} = \frac{Rb}{\sqrt{b^2 - a^2}} \intx\int_{ }^{ } \frac{\tecxt\text{d}z}{\sqrt{1 - \kla\left( \frac{bRz}{\sqrt{b^2 - a^2}} \right)^2}} = \\
&= \intx\int_{ }^{ } \frac{ \mathrm{d}z'}{\sqrt{1 - z'^2}} = \arcsin z' = \arcsin \frac{Rbz}{\sqrt{b^2 - a^2}}.
\end{align}
Thus the full solution is given by
\begin{align}
R\cos\theta(t) &= z(t) = R\sqrt{1 - \frac{a^2}{b^2}} \sin(bt +  c).
\end{align}
We may choose our coordinates such that $\phi_0 = 0 = \phi_{\tecxt\text{end}}$ and $\theta_0=0$ corresponds to our starting position.
We may also choose $\dot{\phi}_0 = 0$ due to the symmetry of the angle $\phi$ at the start and at the end.
Then $a\equiv 0$ and we find the very simple trajectory
\begin{align}
z(t) &= R\sin (bt + c) \equiv R\cos (\dot{\theta}_0 t).
\end{align}
This is equivalent to $\theta(t) = \dot{\theta}_0 t$ where $\phi \equiv\tecxt\text{const}$.
Thus the geodesics on a sphere between two points $A$ and $B$ are always given by the intersection of the surface of the sphere and the plane containing the origin $O$ as well as the points $A$ and $B$.

\section{Christoffel symbols and Ricci-scalar for special 4D coordinates}
We are given the space-time interval $\tecxt\text{d}t^2 = f(r) \tecxt\text{d}r^2 + r^2\tecxt\text{d}\Omega^2 + h(r)\tecxt\text{d}\tau^2$ for spherical 4D coordinates.
Then the Lagrangian is given by
\begin{align}
L &= g_{\mu\nu}\dot{x}^\mu \dot{x}^\nu = f\dot{r}^2 + r^2\dot{\theta}^2 + r^2\sin^2\theta \dot{\phi}^2 + h\dot{\tau}^2.
\end{align}
The Euler-Lagrange equations are
\begin{align}
0 &= \frac{\tecxt\text{d}}{\tecxt\text{d}t} \frac{\partial L}{\partial \dot{r}} - \frac{\partial L}{\partial r} = \frac{\tecxt\text{d}}{\tecxt\text{d}t} \kla\left( 2f\dot{r} \right) - f'\dot{r}^2 - 2r\dot{\theta}^2 - 2r\sin^2\theta \dot{\phi}^2 - h'\dot{\tau}^2 = \\
&= 2f\kla\left( \ddot{r} + \frac{f'\dot{r}^2}{2f} - \frac{r}{f}\dot{\theta}^2 - \frac{r\sin^2\theta}{f} \dot{\phi}^2 - \frac{h'}{2f} \dot{\tau}^2 \right),\\
0 &= \frac{\tecxt\text{d}}{\tecxt\text{d}t} \frac{\partial L}{\partial \dot{\theta}} - \frac{\partial L}{\partial \theta} =  \frac{\tecxt\text{d}}{\tecxt\text{d}t} \kla\left( 2r^2\dot{\theta} \right) - 2r^2\sin\theta\cos\theta \dot{\phi}^2 = \\
&= 2r^2 \kla\left( \ddot{\theta} + \frac{2}{r}\dot{r}\dot{\theta} - \sin\theta\cos\theta \dot{\phi}^2 \right),\\
0 &= \frac{\tecxt\text{d}}{\tecxt\text{d}t} \frac{\partial L}{\partial \dot{\phi}} - \frac{\partial L}{\partial \phi} =  \frac{\tecxt\text{d}}{\tecxt\text{d}t} \kla\left( 2r^2\sin^2\theta \dot{\phi} \right) =  2r^2\sin^2\theta \kla\left( \ddot{\phi} + \frac{2}{r}\dot{r}\dot{\phi} + \frac{2\sin\theta\cos\theta}{\sin^2\theta} \dot{\theta}\dot{\phi} \right),\\
0 &= \frac{\tecxt\text{d}}{\tecxt\text{d}t} \frac{\partial L}{\partial \dot{\tau}} - \frac{\partial L}{\partial \tau} = \frac{\tecxt\text{d}}{\tecxt\text{d}t} \kla\left( 2h\dot{\tau} \right) = 2h\kla\left( \ddot{\tau} + \frac{h'}{h} \dot{r}\dot{\tau} \right) .
\end{align}
Thus we can read of all non-zero Christoffel symbols
\begin{align}
\Gamma^r_{rr} &= \frac{f'}{2f}, &&& \Gamma^r_{\theta\theta} &= -\frac{r}{f}, &&& \Gamma^r_{\phi\phi} &= -\frac{r\sin^2\theta}{f}, &&& \Gamma^r_{\tau\tau} &= -\frac{h'}{2f}, \\
\Gamma^\theta_{r\theta} &= \frac{1}{r}, &&& \Gamma^\theta_{\phi\phi} &= -\sin\theta\cos\theta,\\
\Gamma^\phi_{r\phi} &= \frac{1}{r}, &&& \Gamma^\phi_{\theta\phi} &= \cot\theta,\\
\Gamma^\tau_{r\tau} &= \frac{h'}{2h}.
\end{align}
In order to compute the Ricci-scalar, we need the following components of the curvature tensor.
The general formula is $R^\rho_{\ \mu\lambda\nu} = \partial_\lambda \Gamma^\rho_{\mu\nu} - \partial_\nu \Gamma^\rho_{\mu\lambda} + \Gamma^\sigma_{\mu\nu} \Gamma^\rho_{\sigma\lambda} - \Gamma^\sigma_{\mu\lambda} \Gamma^\rho_{\sigma\nu}$.
For a diagonal metric, we may also write $R_{\rho\mu\lambda\nu} = g_{\rho\sigma} R^\sigma_{\ \mu\lambda\nu} \equiv g_{\rho\rho} R^\rho_{\ \mu\lambda\nu}$
\begin{align}
R_{r\theta r\theta} &= g_{rr} R^r_{\ \theta r\theta} = f \kla\left( \partial_r \Gamma^r_{\theta\theta} - \partial_\theta \Gamma^r_{\theta r} + \Gamma^\sigma_{\theta\theta}\Gamma^r_{\sigma r} - \Gamma^\sigma_{\theta r} \Gamma^r_{\sigma\theta} \right) = \\
&= f\kla\left( -\frac{1}{f} + \frac{rf'}{f^2} + \frac{-r}{f} \frac{f'}{2f} - \frac{1}{r} \frac{-r}{f} \right) = \frac{rf'}{f} - \frac{rf'}{2f} = \frac{rf'}{2f},\\
R_{r\phi r\phi} &= g_{rr} R^r_{\ \phi r\phi} = f \kla\left( \partial_r \Gamma^r_{\phi\phi} - \partial_\phi \Gamma^r_{\phi r} + \Gamma^\sigma_{\phi\phi}\Gamma^r_{\sigma r} - \Gamma^\sigma_{\phi r} \Gamma^r_{\sigma\phi} \right) = \\
&= f \kla\left( -\frac{\sin^2\theta}{f} + \frac{rf'\sin^2\theta}{f^2} + \frac{-r\sin^2\theta}{f}\frac{f'}{2f} - \frac{1}{r} \frac{-r\sin^2\theta}{f} \right) = \sin^2\theta R_{r\theta r\theta},\\
R_{r\tau r\tau} &= g_{rr} R^r_{\ \tau r\tau} = f \kla\left( \partial_r \Gamma^r_{\tau\tau} - \partial_\tau \Gamma^r_{\tau r} + \Gamma^\sigma_{\tau\tau}\Gamma^r_{\sigma r} - \Gamma^\sigma_{\tau r} \Gamma^r_{\sigma\tau} \right) = \\
&= f \kla\left( \frac{-h''}{2f} + \frac{h'f'}{2f^2} + \frac{-h'}{2f} \frac{f'}{2f} - \frac{h'}{2h} \frac{-h'}{2f} \right) = \frac{h'f'}{4f} - \frac{h''}{2} + \frac{(h')^2}{4h},\\
R_{\theta\phi \theta\phi} &= g_{\theta\theta} R^\theta_{\ \phi \theta\phi} = r^2 \kla\left( \partial_\theta \Gamma^\theta_{\phi\phi} - \partial_\phi \Gamma^\theta_{\phi \theta} + \Gamma^\sigma_{\phi\phi}\Gamma^\theta_{\sigma \theta} - \Gamma^\sigma_{\phi \theta} \Gamma^\theta_{\sigma\phi} \right) = \\
&= r^2 \kla\left( \sin^2\theta - \cos^2\theta + \frac{-r\sin^2\theta}{f} \frac{1}{r} - \cot\theta (-\sin\theta\cos\theta) \right) = \\
&= r^2\sin^2\theta \kla\left( 1 - \frac{1}{f} \right),\\
R_{\theta\tau \theta\tau} &= g_{\theta\theta} R^\theta_{\ \tau \theta\tau} = r^2 \kla\left( \partial_\theta \Gamma^\theta_{\tau\tau} - \partial_\tau \Gamma^\theta_{\tau \theta} + \Gamma^\sigma_{\tau\tau}\Gamma^\theta_{\sigma \theta} - \Gamma^\sigma_{\tau \theta} \Gamma^\theta_{\sigma\tau} \right) = \\
&= r^2 \kla\left( \frac{-h'}{2f} \frac{1}{r} \right) = -\frac{rh'}{2f},\\
R_{\phi\tau \phi\tau} &= g_{\phi\phi} R^\phi_{\ \tau \phi\tau} = r^2\sin^2\theta \kla\left( \partial_\phi \Gamma^\phi_{\tau\tau} - \partial_\tau \Gamma^\phi_{\tau \phi} + \Gamma^\sigma_{\tau\tau}\Gamma^\phi_{\sigma \phi} - \Gamma^\sigma_{\tau \phi} \Gamma^\phi_{\sigma\tau} \right) = \\
&= r^2\sin^2\theta \kla\left( \frac{-h'}{2f} \frac{1}{r} \right) = \sin^2\theta R_{\theta\tau \theta\tau}.
\end{align}
In order to compute the Ricci-scalar from these components, we make use of an identity valid for diagonal metrics
\begin{align}
R &= g^{\mu\nu} R_{\mu\nu} = g^{\mu\nu} g^{\rho\sigma} R_{\sigma\mu \rho\nu} \equiv \sum_{\mu\rho} g^{\mu\mu} g^{\rho\rho} R_{\rho\mu\rho\mu} = 2\sum_{\mu > \rho} g^{\mu\mu} g^{\rho\rho} R_{\rho\mu\rho\mu}.
\end{align}
Thus we have
\begin{align*}
R &= \frac{2}{f} \kla\left( \frac{1}{r^2} R_{r\theta r\theta} + \frac{1}{r^2\sin^2\theta} R_{r\phi r\phi} + \frac{1}{h} R_{r\tau r\tau} \right) + \\
&\hspace{0.5cm} + \frac{2}{r^2} \kla\left( \frac{1}{r^2\sin^2\theta} R_{\theta\phi \theta\phi} + \frac{1}{h} R_{\theta\tau \theta\tau} \right) + \frac{2}{r^2\sin^2\theta} \frac{1}{h} R_{\phi\tau \phi\tau} = \numberthis  \\
&= \frac{1}{r^2}\frac{rf'}{f^2} + \frac{\sin^2\theta}{r^2\sin^2\theta} \frac{rf'}{f^2} + \frac{1}{fh} \kla\left( \frac{h'f'}{2f} - h'' + \frac{(h')^2}{2h} \right) + \\
&\hspace{0.5cm} + \frac{2}{r^2} \kla\left( 1 - \frac{1}{f} \right) + \frac{2}{r^2h} \frac{-rh'}{2f} + \frac{1}{rh} \frac{-h'}{f} =  \numberthis \\
&= \frac{2f'}{rf^2} + \frac{f'h'}{2f^2h} - \frac{h''}{fh} + \frac{(h')^2}{2fh^2} + \frac{2}{r^2} \kla\left( 1 - \frac{1}{f} \right) - \frac{2h'}{rfh}. \numberthis 
\end{align*}

\section{Curvature of an ellipsoid}
Consider an ellipsoid parametrized by 
\begin{align}
\vec{r} &= R\begin{pmatrix}
\cos\phi\sin\theta \\ \sin\phi\sin\theta \\ \alpha\cos\theta
\end{pmatrix}
\end{align}
with $R,\alpha = \tecxt\text{const}$.
The line element is then given by
\begin{align}
\tecxt\text{d}\vec{r} &= R\begin{pmatrix}
-\sin\phi\sin\theta \tecxt\text{d}\phi + \cos\phi\cos\theta \tecxt\text{d}\theta \\ 
\cos\phi\sin\theta \tecxt\text{d}\phi + \sin\phi\cos\theta \tecxt\text{d}\theta \\
-\alpha \sin\theta \tecxt\text{d}\theta
\end{pmatrix},
\end{align}
which squares to 
\begin{align}
(\tecxt\text{d}\vec{r})^2 &= R^2 \kla\left( \sin^2\theta \tecxt\text{d}\phi^2 + \cos^2\theta \tecxt\text{d}\theta^2 + \alpha^2\sin^2\theta \tecxt\text{d}\theta^2 \right) = \\
&= R^2 \kla\left( \klam\left[ 1  + (\alpha^2 - 1)\sin^2\theta \right] \tecxt\text{d}\theta^2 + \sin^2\theta \tecxt\text{d}\phi^2 \right)•
\end{align}
The metric results in
\begin{align}
(g_{ab}) &= R^2 \begin{pmatrix}
1 + (\alpha^2 - 1)\sin^2\theta & 0 \\ 0 & \sin^2\theta
\end{pmatrix}.
\end{align}
Let $\beta := \alpha^2 - 1$.
Then the Lagrangian is
\begin{align}
L &= g_{ab} \dot{x}^a \dot{x}^b = R^2 \kla\left( [1 + \beta \sin^2\theta] \dot{\theta}^2 + \sin^2\theta\dot{\phi}^2 \right).
\end{align}
The Euler-Lagrange equations become
\begin{align}
0 &= \frac{\tecxt\text{d}}{\tecxt\text{d}t} \frac{\partial L }{\partial \dot{\theta}} - \frac{\partial L}{\partial \theta} \propto \\
&\propto \frac{\tecxt\text{d}}{\tecxt\text{d}t} \kla\left( 2[1 + \beta \sin^2\theta] \dot{\theta} \right) - 2\beta \sin\theta\cos\theta \dot{\theta}^2 - 2\sin\theta\cos\theta \dot{\phi}^2 = \\
&= 2[1 + \beta \sin^2\theta] \ddot{\theta} + 2\beta \sin\theta\cos\theta \dot{\theta}^2 - 2\sin\theta\cos\theta \dot{\phi}^2,\\
0 &= \frac{\tecxt\text{d}}{\tecxt\text{d}t} \frac{\partial L }{\partial \dot{\phi}} - \frac{\partial L}{\partial \phi} \propto \\
&\propto \frac{\tecxt\text{d}}{\tecxt\text{d}t} \kla\left( 2\sin^2\theta \dot{\phi} \right) = 2\sin^2\theta \ddot{\phi} + 4\sin\theta \cos\theta \dot{\theta} \dot{\phi}.
\end{align}
Thus the Christoffel symbols are
\begin{align}
\Gamma^\theta_{\theta\theta} &= \frac{\beta \sin\theta\cos\theta}{1 + \beta \sin^2\theta}, &&& \Gamma^\theta_{\phi\phi} &= -\frac{\sin\theta \cos\theta}{1 + \beta \sin^2\theta},\\
\Gamma^\phi_{\theta\phi} &= \cot\theta.
\end{align}
The curvature tensor in two dimensions only has one independent component
\begin{align*}
R_{\theta\phi\theta\phi} &= g_{\theta \sigma} R^\sigma_{\ \phi\theta \phi} = g_{\theta\theta} \kla\left( \partial_\theta \Gamma^\theta_{\phi\phi} - \partial_\phi \Gamma^\theta_{\phi\theta} + \Gamma^\sigma_{\phi\phi} \Gamma^\theta_{\sigma\theta} - \Gamma^\sigma_{\phi\theta} \Gamma^\theta_{\sigma\phi} \right) = \numberthis \\
&= R^2 [1 + \beta\sin^2\theta] \kla\bigg( \frac{\sin^2\theta - \cos^2\theta}{1 + \beta\sin^2\theta} + \frac{2\beta\sin^2\theta\cos^2\theta }{(1 + \beta\sin^2\theta)^2} + \\
&\hspace{3.5cm} + \frac{-\sin\theta\cos\theta}{ 1 + \beta\sin^2\theta } \frac{\beta\sin\theta\cos\theta }{1 + \beta\sin^2\theta} - \cot\theta \frac{-\sin\theta\cos\theta }{1 + \beta\sin^2\theta} \bigg) = \numberthis \\ 
&= R^2\sin^2\theta + R^2\frac{\beta \sin^2\theta\cos^2\theta }{1 + \beta\sin^2\theta} = R^2\sin^2\theta \kla\left( 1 + \frac{\beta\cos^2\theta}{1 + \beta\sin^2\theta} \right) = \\
&= R^2\sin^2\theta \frac{1 + \beta}{1 + \beta\sin^2\theta}.
\end{align*}
Note that the components of the Riemann tensor then become
\begin{align}
R^\theta_{\ \phi\theta\phi} &= -R^\theta_{\ \phi\phi\theta} = g^{\theta\theta} R_{\theta\phi \theta\phi} = \sin^2\theta \frac{1 + \beta}{(1 + \beta\sin^2\theta)^2},\\
R^\phi_{\ \theta\phi\theta} &= -R^\phi_{\ \theta\theta\phi} = g^{\phi\phi} R_{\theta\phi\theta\phi} = \frac{1 + \beta}{1 + \beta\sin^2\theta}.
\end{align}
The components of the Ricci-tensor are then given by
\begin{align}
R_{\theta\theta} &= R^\sigma_{\ \theta\sigma\theta} = \frac{1 + \beta}{1 + \beta\sin^2\theta},\\
R_{\phi\phi} &= R^\sigma_{\ \phi\sigma\phi} = \sin^2\theta \frac{1 + \beta}{(1 + \beta\sin^2\theta)^2}.
\end{align}
Then the Ricci-scalar can be computed
\begin{align}
R &= g^{ab}R_{ab} = g^{\theta\theta} R_{\theta\theta} + g^{\phi\phi} R_{\phi\phi} = \\
&=  \frac{2}{R^2}\frac{1 + \beta}{(1 + \beta\sin^2\theta)^2} = \frac{2}{R^2} \frac{\alpha^2}{\klam\left[ 1 + (\alpha^2 - 1)\sin^2\theta \right]^2}.
\end{align}

%\section{Curvature of a sphere}
%The Christoffel symbols are
%\begin{align}
%\Gamma^\theta_{\phi\phi} &= -\sin\theta \cos\theta,\\
%\Gamma^\phi_{\theta\phi} &= \cot\theta.
%\end{align}
%The curvature tensor in two dimensions only has one independent component
%\begin{align}
%R_{\theta\phi\theta\phi} &= g_{\theta \sigma} R^\sigma_{\ \phi\theta \phi} = g_{\theta\theta} \kla\left( \partial_\theta \Gamma^\theta_{\phi\phi} - \partial_\phi \Gamma^\theta_{\phi\theta} + \Gamma^\sigma_{\phi\phi} \Gamma^\theta_{\sigma\theta} - \Gamma^\sigma_{\phi\theta} \Gamma^\theta_{\sigma\phi} \right) = \\
%&= R^2 \kla\left( \sin^2\theta - \cos^2\theta  - \cot\theta (-\sin\theta\cos\theta)  \right) = R^2\sin^2\theta
%\end{align}
%Note that the components of the Riemann tensor then become
%\begin{align}
%R^\theta_{\ \phi\theta\phi} &= -R^\theta_{\ \phi\phi\theta} = g^{\theta\theta} R_{\theta\phi \theta\phi} = \sin^2\theta,\\
%R^\phi_{\ \theta\phi\theta} &= -R^\phi_{\ \theta\theta\phi} = g^{\phi\phi} R_{\theta\phi\theta\phi} = 1.
%\end{align}
%The components of the Ricci-tensor are then given by
%\begin{align}
%R_{\theta\theta} &= R^\sigma_{\ \theta\sigma\theta} = 1,\\
%R_{\phi\phi} &= R^\sigma_{\ \phi\sigma\phi} = \sin^2\theta.
%\end{align}
%Then the Ricci-scalar can be computed
%\begin{align}
%R &= g^{ab}R_{ab} = g^{\theta\theta} R_{\theta\theta} + g^{\phi\phi} R_{\phi\phi} =  \frac{2}{R^2}.
%\end{align}

%\begin{thebibliography}{9}
%\bibitem{example} description, \url{https://www.exmapleurl}, day.\,month~2021.
%\end{thebibliography}

\end{document}